\section{Derivadas regularizadas y valores espectrales impares}
\label{ix-add:bendersky}

La continuación y la ecuación funcional demostradas en la
proposición~\ref{rz-num-03-2} permiten relacionar dos operaciones
sobre el mismo funcional $Z$: la derivación en enteros negativos y la
evaluación convergente en enteros positivos impares. Las coordenadas
regionales y el operador de numeración mantienen su genealogía anterior.
La notación $\pi$ designa la coordenada HMT que ya intervino en la
construcción theta--Mellin.

Sean $H_k=\sum_{j=1}^k1/j$, con $H_0=0$, y $\mathsf B_j$ los
números de Bernoulli construidos por la recurrencia de la
sección~\ref{lectura:manuscrito-04}, con $\mathsf B_1=-1/2$.
La tipografía $\mathsf B_j$ distingue estos coeficientes de la constante
de Meissel--Mertens. Se define la jerarquía de Bendersky--Glaisher--Kinkelin
mediante
\[
 A_k=\exp\!\left(\frac{H_k\mathsf B_{k+1}}{k+1}-Z'(-k)\right),
 \qquad k\in\mathbb N_0.
\]
Los puntos $-k$ son regulares. El término de Bernoulli fija la
normalización correspondiente a cada nivel antes de su evaluación.

\subsection{Covariancia de las derivadas regularizadas}
Se utiliza aquí el funcional $Z_\lambda=\lambda^{-s}Z$, cuyo residuo
es $\lambda^{-1}$, y se define
\[
 D_k(\lambda)=\exp[-Z_\lambda'(-k)].
\]
Este objeto conserva la normalización del operador escalado, distinta
de la normalización de residuo unitario empleada para
$\widehat Z_\lambda$ en la sección~\ref{ix-add:stieltjes}.

\HMTresultado{Proposición}{Transporte de las derivadas regularizadas}{ix-add:regularized-scaling}
Para $\lambda>0$ y $k\in\mathbb N_0$,
\[
 Z_\lambda'(-k)=\lambda^k
       \bigl(Z'(-k)-(\log\lambda)Z(-k)\bigr),
\]
\[
 D_k(\lambda)=D_k(1)^{\lambda^k}
               \lambda^{\lambda^k Z(-k)},\qquad
 D_k(1)=A_k\exp\!\left(-\frac{H_k\mathsf B_{k+1}}{k+1}\right).
\]
\textbf{Demostración.} La regla del producto aplicada a
$\lambda^{-s}Z(s)$ da la primera identidad. Su exponencial negativa
da la segunda, y la definición de $A_k$ da la tercera.
\hfill$\square$

En los niveles iniciales,
\[
\begin{aligned}
 A_0&=e^{-Z'(0)},\\
 \log A_1&=\frac1{12}-Z'(-1),\qquad \log A_2=-Z'(-2),\\
 \log A_3&=-\frac{11}{720}-Z'(-3),\qquad \log A_4=-Z'(-4).
\end{aligned}
\]
Por ejemplo, $D_1(1)=A_1e^{-1/12}$. Mantener el normalizador evita
identificar dos lectores diferentes por compartir la misma derivada.

\Needspace{12\baselineskip}
\subsection{Reconstrucción de los valores impares}
\HMTresultado{Teorema}{Correspondencia entre niveles negativos pares y positivos impares}{ix-add:odd-values}
Para todo entero $n\geq1$,
\[
 Z'(-2n)=(-1)^n\frac{(2n)!}{2(2\pi)^{2n}}Z(2n+1),
\]
\[
 \frac{\log D_{2n}(\lambda)}{\lambda^{2n}}
 =\log A_{2n}
 =(-1)^{n+1}\frac{(2n)!}{2(2\pi)^{2n}}Z(2n+1).
\]
\textbf{Demostración.} La ecuación funcional de la forma completada
implica la identidad meromorfa
\[
 \pi^{-s/2}\Gamma(s/2)Z(s)
 =\pi^{-(1-s)/2}\Gamma((1-s)/2)Z(1-s).
\]
En $s=-2n$, el segundo miembro es regular y el factor gamma del
primero tiene un polo simple. Por tanto $Z(-2n)=0$. La recurrencia
de gamma y su residuo en cero dan
\[
 \Gamma(s/2)=\frac{2(-1)^n}{n!(s+2n)}+O(1).
\]
Además,
\[
 \Gamma\!\left(n+\frac12\right)
 =\frac{(2n)!\sqrt\pi}{4^nn!}.
\]
Comparar los términos constantes en la identidad meromorfa produce
\[
 \pi^n\frac{2(-1)^n}{n!}Z'(-2n)
 =\pi^{-n}\frac{(2n)!}{4^nn!}Z(2n+1),
\]
que es la primera fórmula. La serie
$t/(e^t-1)+t/2$ es par; por ello $\mathsf B_{2n+1}=0$ para
$n\geq1$. La definición de $A_{2n}$ y la proposición anterior
completan la prueba. \hfill$\square$

En particular, el primer nivel par positivo satisface
\[
 Z(3)=4\pi^2\log A_2.
\]
La constante de Apéry evaluada con resto en la
sección~\ref{lectura:manuscrito-04} queda así relacionada con una
derivada regularizada del mismo funcional. Los intervalos allí
demostrados se transportan por la exponencial monótona una vez
conservadas las cotas de la coordenada regional $\pi$.

El nivel cero tiene una evaluación distinta. La ecuación funcional,
junto con reflexión y duplicación de gamma, equivale a
\[
 Z(s)=2^s\pi^{s-1}\sin(\pi s/2)\Gamma(1-s)Z(1-s).
\]
La fórmula de Euler
\[
 \Gamma(1+z)=\lim_{m\to\infty}
 m^z\prod_{j=1}^{m}(1+z/j)^{-1}
\]
da $\log\Gamma(1+z)=-\gamma_0z+O(z^2)$: el coeficiente lineal es
el límite de $\log m-H_m$, y el resto cuadrático queda controlado
por la serie convergente de $j^{-2}$. Usando
$Z(1-s)=-1/s+\gamma_0+O(s)$ y
$\Gamma(1-s)=1+\gamma_0s+O(s^2)$, los términos de grado uno
que contienen $\gamma_0$ se cancelan y resulta
\[
 Z(s)=-\frac12-\frac12\log(2\pi)s+O(s^2).
\]
En consecuencia,
\[
 A_0=\sqrt{2\pi},\qquad D_0(\lambda)=\sqrt{\frac{2\pi}{\lambda}}.
\]
La condición $n\geq1$ del teorema conserva esta diferencia: el nivel
cero no se obtiene sustituyendo $Z(1)$, que tiene un polo, en la fórmula
de valores impares. Las identidades de gamma expresan la normalización
analítica posterior, sin intervenir en la selección de las semillas.
